\section{The closed source proof of \texorpdfstring{$\mathrm{OneCharU\_s}$}{OneCharU\_s}}
\label{sec:kt-onechar-closed}

\paragraph{Finding KT-ONECHAR-CLOSED.}
Under the stipulated convention that the archive declarations are cited as
theorems, \path{UpperTheory.V.Z.J.OneCharU_s} has a closed proof term with
no additional mathematical hypotheses.  This is a direct source trace,
not a new Lean kernel build.  No Lean executable was available in this
session, and the accompanying file
\path{OneCharUClosed_UNCOMPILED.lean} is explicitly marked uncompiled.
The universal arithmetic conditions inside the definition of
\path{OneCharU_s}, including \path{GD}, remain part of the statement.

All locators in this section use the pinned commit
\path{7cbf7a0c32fca1badf05931879a37c3f8a2b03b4} of
\url{https://github.com/CocoSgt/riemann-hypothesis-lean}.
Abbreviate
\path{src/Development/AutomorphicArithmetic/AutomorphicVarieties/}
by AV/ and
\path{AV/ModularityArgument/ThresholdArguments/DenseThresholdCases/}
by DTC/.  The actual implication is
\path{UpperTheory.V.Z.J.oneCharU_of_katoTwoW_rig_admI_automorphy},
DTC/\path{FOneCharU.lean}:75--95.  Its nine hypotheses, in order, are
the following; all namespace prefixes shown are exact.

\begin{enumerate}
\item \path{UpperTheory.V.U3.CohInp}, supplied by
\path{UpperTheory.V.HV.Core0.cohInp_alt hH.litN hK},
DTC/\path{CohInp.lean}:16--17.
\item \path{UpperTheory.V.Alg.LemmaM}, supplied by
\path{UpperTheory.V.Alg.M.main_holds},
AV/\path{CommutativeAlgebra/GradedModulesLength/MainCotSpan.lean}:32--70.
\item \path{UpperTheory.V.Alg.LemmaM'}, supplied by
\path{UpperTheory.V.Alg.M.mprime_holds},
AV/\path{CommutativeAlgebra/GradedModulesLength/Mprime.lean}:38--80.
\item \path{UpperTheory.V.Alg.LemmaP}, supplied by
\path{UpperTheory.V.Alg.lemmaP},
AV/\path{CommutativeAlgebra/LemmaP.lean}:256--424.
\item \path{UpperTheory.V.Z.Kato.WeilSphericalRecords.KatoTwoW},
supplied by
\path{UpperTheory.V.Z.Kato.WeilSphericalRecords.katoTwoW_of_printsH hH},
DTC/\path{KatoTwoWOfPublishedH.lean}:48--49.
\item \path{UpperTheory.V.Th.RigidityFirstCase}, supplied by
\path{UpperTheory.V.Z.J.rigIUE_alt hC hH hL hE},
DTC/\path{ERigIU.lean}:80--146.
\item \path{UpperTheory.V.Th.RigiditySecondCase}, supplied by
\path{UpperTheory.V.Z.Win.RR.rigII_a_alt hC hH hL},
DTC/\path{RRRigII.lean}:25--102.
\item \path{UpperTheory.V.Z.J.AdmI_s}, supplied by
\path{UpperTheory.V.Z.J.admissible_typeI_even_level_alt hC hH hE},
DTC/\path{AdmissibleTypeIEvenLevel.lean}:76--136.
\item \path{UpperTheory.V.Aut.AutomorphyUnramifiedCoefficientsAlt},
supplied by
\path{UpperTheory.V.HV.Core0.automorphyUnramified_of_litN_of_gdEnd}
applied to \path{hH.litN} and
\path{(UpperTheory.V.Z.Win.Dn.ends0_alt hH).1};
DTC/\path{CohInp.lean}:18 and DTC/\path{DnEndsTop.lean}:15--16.
\end{enumerate}

Here \path{hC} is the term in item 1;
\path{hK := UpperTheory.V.ZK.lemmaK_litN_alt hH.litN}
uses DTC/\path{KFinal.lean}:11;
\path{hL := UpperTheory.V.HV.DescentVanishingReductions.adjoint_lifting_of_cohomology hH}
uses DTC/\path{AdjointLiftingOfCohomology.lean}:31; and
\path{hE := unlEnd3E (UpperTheory.V.Z.J.linkedUnlE_alt hH)}.
The latter application is DTC/\path{LinkedUnlE.lean}:51;
\path{linkedUnlE_alt} is declared there at lines 43--48 and
\path{unlEnd3E} at
AV/\path{ModularityArgument/LinkedUnramifiedChains/LogarithmicNormBounds/BlockValuationRep.lean}:234.

\paragraph{Proof of the closedness claim.}
The application with all nine proof terms is written literally in
DTC/\path{LinkedUnlE.lean}:49--54, in the proof of
\path{count_i_of_printed_bound_k_kato_automorphy}.
Its caller, DTC/\path{WindowBoundOfLiteratureLevel.lean}:25--29,
supplies \path{hK}, the Kato term in item 5, and the automorphy term in
item 9.  The remaining record can be closed by setting
\path{hH} equal to
\path{published_hypotheses_from_patching_alt}
applied to \path{UpperTheory.V.SmallClassNumbers.litN}
and \path{UpperTheory.V.Z.Tf.SR.descWeak}.
The record constructor is
DTC/\path{PublishedHypothesesFromPatching.lean}:98--103,
and \path{SmallClassNumbers.litN} is the closed theorem in
AV/\path{SmallClassNumbers/CorGKOf.lean}:20--23.
The other argument, \path{UpperTheory.V.Z.Tf.SR.descWeak}, is
the closed theorem at
AV/\path{ModularityArgument/TraceFormula/SplitRigidity/HasOfSplit.lean}:29--30,
from \path{descWeakAll} at lines 26--27.
Indeed, \path{CorGKOf.lean}:10--18 repeats the same inner
\path{OneCharU_s} proof term, and line 25 supplies precisely this closed
record to prove \path{corGK}.  Taking that inner term after the same
substitution proves the claim.  No inference from mere membership in
a syntactic dependency closure is needed. \hfill$\square$

\paragraph{The actual route from the named root.}
\path{src/Main.lean}:4 defines \path{riemannHypothesis := rhOfAll}.
\path{src/Supplements/Kernels/MainTheorems.lean}:7--9 takes
\path{rhOfAll} from \path{allChi}, itself
\path{UpperTheory.FinalTheorems.everyCharacter}.
\path{src/Development/AutomorphicArithmetic/CharacterBoxes/MainConclusion.lean}:7--11
applies \path{everyCharacter_of_row} to
\path{UpperTheory.Gal.Int.galoisLocalGlobalStrong_recParameter}
and uses \path{UpperTheory.One.vwin_of_row}.
DTC/\path{VwinOfRow.lean}:4--5 supplies the open-image record,
\path{UpperTheory.BaseChangeIdentities.baseChangeN}, and that same
Galois row to \path{automorphic_window_of_inputs}.
DTC/\path{TopMainZj.lean}:71--74 and 55--58 construct the required
\path{LiteratureRankNHypothesis} record.
DTC/\path{WindowBoundOfLiteratureLevel.lean}:30--34 and 25--29 then
reach the literal nine-term application already identified.
The Galois-row theorem itself is closed at
\path{src/Development/AutomorphicArithmetic/Foundations/GaloisRepresentations/IntegralStructures/RowSrc.lean}:6.
These declarations are cited archive inputs; their entire transitive
dependency tree was not re-proved in this session.

\paragraph{Exact indexing and multiplicity.}
DTC/\path{NumberU3US.lean}:57--59 makes $X$ a finite set of actual
three-dimensional Galois representations.  Besides admissibility it
requires
\[
\rho,\rho'\in X,\quad\operatorname{RepIso}(\rho,\rho')
\quad\Longrightarrow\quad\rho=\rho'.
\]
This is distinctness of isomorphism classes, not distinctness of one
Frobenius trace.  DTC/\path{FOneCharU.lean}:59--65 proves the
injectivity required to pass to $X$; lines 66--74 use
\path{Finset.sum_image}.  Thus a unit criterion needs pointwise coverage,
whereas a numerical product bound must keep enough eigenform indices
to cover all pairwise nonisomorphic systems, even if some traces coincide.

\paragraph{The range of the rational specialization.}
AV/\path{SmallClassNumbers/ClassNumberThree.lean}:64--86 proves
\path{GD} for the trivial character modulo $2$, $K=\mathbf Q(\sqrt{-3})$,
every even $e\ge1$, and every prime $\ell\ge5$ satisfying
$(\ell-1)\nmid2e$.  Its statement contains no inequality $e<\ell$.
The exclusion $\ell\ne7$ belongs to the subsequent Frobenius-$7$
evaluation, not to this \path{GD} theorem.

\paragraph{Interpretation and verification.}
This closes the source-standing gate in K17-S17-UNIT.  It supplies the
archive's bound on strict Galois-cohomology extension classes of Tate
twists over $\operatorname{Spec}\mathbf Z[1/\ell]$.
Vanishing still requires the arithmetic trace input and the carried
K17-VK-PT comparison.  The complete extraction proof and its source
applications were read and recorded; no new kernel run or global
axiom audit is asserted.
